% Volume VI: Adversaries, Platforms, and Automated Judgment
% Part XI. Untrustworthy Actors
% Chapter 60: Trust Networks and Correlated Capture
% Spine: proof-spines/ch060-spine.md

\chapter{Trust Networks and Correlated Capture}
\label{ch:ch060}


Trust should be evaluated over dependency graphs rather than by counting repetitions. The key danger is \textbf{correlated capture}: what looks like many confirming voices may in fact be many downstream copies of one compromised source. If reports \(R_1,\ldots,R_n\) share the same ancestor \(S\), their agreement does not supply the evidentiary force that genuine independence would have supplied.
\ToyModel\ The shared-ancestor sketch models apparent plurality as hidden dependence.

The outline's point is therefore structural rather than psychological. A thousand repetitions can remain only one underlying observation, one leak, one institution, or one gatekeeper echoed through many channels. Their joint evidentiary weight is then far smaller than it would be under independence, even if no single downstream speaker is consciously deceptive. Resilient trust networks must accordingly diversify origins, incentives, and review paths, because nominal plurality is useless wherever hidden dependence has already collapsed into \emph{correlated capture}.
